\documentclass[10pt,a4paper]{amsart}
\usepackage[margin=19mm]{geometry}
\usepackage{amsmath,amssymb,mathtools}
\usepackage[scr=rsfs,cal=euler]{mathalfa}
\usepackage{xfrac}
\usepackage[unicode,hidelinks,bookmarks=false]{hyperref}
\newcommand{\ie}{i.e.\ }
\newcommand{\cf}{cf.\ }
\newcommand{\resp}{respectively\ }
\newcommand{\Subsection}{\subsection}
\providecommand{\nobreakdash}{\nobreak\hbox{-}}
\setlength{\emergencystretch}{2em}
\multlinegap0pt
\usepackage{xcolor}
\usepackage{algorithm}\usepackage{algpseudocode}
\usepackage{amssymb}
\usepackage{bm}
\usepackage{float}
\usepackage{tikz}
\newtheorem{theorem}{Theorem}
\newcommand{\C}{\mathbb{C}}
\newcommand{\RP}{\mathbb{R}\mathbb{P}}
\newcommand{\p}{{\mathbb{P}}}
\title{Introduction to non-Abelian Patchworking}
\author{Turgay Akyar, Mikhail Shkolnikov}
\date{Source: arXiv:2603.08094v1 (CC BY 4.0)}
\begin{document}
\maketitle

\noindent\emph{Source and licence.} Introduction to non-Abelian Patchworking. Version arXiv:2603.08094v1 (CC BY 4.0), distributed by arXiv under the Creative Commons Attribution 4.0 International licence, http://creativecommons.org/licenses/by/4.0/, as recorded in the arXiv licence metadata for this exact version and shown on the abstract page https://arxiv.org/abs/2603.08094v1. Full licence notice: \texttt{licenses/arxiv-2603.08094-CC-BY-4.0.txt}.\par\smallskip

\noindent\textit{Editorial scope.} Counted unit: the article's theorem theorem\_ec (source0.tex lines 233--247). The article's complete original proof of it is delivered with it (source0.tex lines 248--290, one \texttt{proof} environment of 43 lines). No mathematics is written, repaired or re-derived: the statement and the proof are the article's own lines, in the article's own notation, and no retained line is substituted or re-flowed.  No further source-local context is needed: every cross-reference of every retained line resolves inside the two delivered spans.  Nothing was omitted from any retained span, so the package carries no omission record.  No retained line contains a citation, so the wrapper carries no bibliography and no bibliography entry.  No retained line contains a figure environment, an image-inclusion command or drawing code, so no image is delivered and the package declares no asset dependency.  Editorial formatting only, in addition: an AMS \texttt{amsart} preamble that restates the article's own declarations and macros used by the retained lines, namely the environments \texttt{theorem} on separate counters, together with the standard packages those lines need.  The retained environments print in this document as Theorem 1 in order of appearance, whereas the article prints the same items with the numbering of its own full document.  The complete native source of the article as distributed in the arXiv e-print for this exact version is bundled unchanged as \texttt{sources/195983/source0.tex}.
\begin{theorem}\label{theorem_ec}
	Let $X\subset\C\p^3$ be a complex smooth surface of degree $d$ and $D$ be the $PGL_2(\mathbb{C})$-phase diagram of degree $d$, obtained with the initial data described above. If $D$ is ${I_{T^2}}$-real and $d=2k+1,$ then the Euler characteristic is bound to
	$$\chi(\overline{D^{I_{T^2}}})\in[1,\sigma(X)]\cap (2\mathbb{Z}+1)$$
	%(-8k^3-12k^2+2k+3)/3
	In the even case and $D$ being real ${I_{T^2}}$-real, where $d=2k$,
	$$\chi(\overline{D^{I_{T^2}}})\in[0,\sigma(X)]\cap 2\mathbb{Z}.$$
	%(-8k^3-12k^2+2k+3)/
	Moreover, for $I_\emptyset$-real $D$ \[  \chi(\overline{D^{I_{\emptyset}}})=\left\{
	\begin{array}{ll}
		0, & \text{if } d \text{ is even} \\
		1, & \text{if } d \text{ is odd.} \\
	\end{array} 
	\right. \]
	Lastly, for $I_{S^2}$-real $D$ $$d\geq\chi(\overline{D^{I_{S^2}}})\geq d+\sigma(X)\text{ with }\chi(\overline{D^{I_{S^2}}})\equiv d \quad (\text{mod 2}).$$
\end{theorem}

\begin{proof}
	Let us start with the calculations for $\chi(\overline{D^{I_{T^2}}})$ in the even degree case. We first observe that between any two consecutive critical levels $\gamma_j$ and $\gamma_{j+1}$ the real locus consists of cylindrical pieces whose base is the real smooth curve 
	\[
	C_{j+1}(\mathbb{R}) \subset \mathbb{RP}^1 \times \mathbb{RP}^1
	\]
	of bidegree $(2j+2,2j+2)$. Since these pieces are topologically cylinders, they do not contribute to the Euler characteristic. Therefore, for the computation of $\chi$, the essential contribution comes from the behavior at the critical levels themselves, together with the additional coamoeba contributions in the odd degree case.
	
	The general principle determining the computation at a critical level $\gamma_{j+1}$ is the following: the union of a critical level with its reflection in the other connected component of $PGL_2(\mathbb{R})$ reconstructs the real quadric (with the intersection of the curves counted twice)
	
	$$Q(\mathbb{R}) \cong \mathbb{RP}^1 \times \mathbb{RP}^1.$$
	
	From this we obtain, for $j>1$, a contribution of the form
	
	$$\chi\!\left(Q(\mathbb{R})\right)
	+\chi\!\left(C_j(\mathbb{R})\right)
	+\chi\!\left(C_{j+1}(\mathbb{R})\right)
	-\chi\!\left(C_j(\mathbb{R}) \cap C_{j+1}(\mathbb{R})\right).$$
	
	The case $j=1$ is exceptional: at the first critical level only the curve $C_1$ appears, and hence no additional contribution arises from reflection.
	Therefore, we get 
	
	$$\chi(X)=\sum_{j=0}^{k-1}\big(\chi(\mathbb{RP}^1\times\mathbb{RP}^1)+\chi(C_j)+\chi(C_{j+1})-\chi(C_j\cap C_{j+1})\big),$$
	where $C_j$ is a bidegree $(2j,2j)$ smooth curve in $\mathbb{RP}^1\times\mathbb{RP}^1$ for each $j>0$, and $C_0$ is the empty set. We may write $0\leq\chi(C_l\cap C_{l+1})\leq 2(2j)(2j+2)$. This leaves 
	$$0\geq \chi(X)\geq -\sum_{l=0}^{k-1}2(2j)(2j+2)=\dfrac{-8k^3+8k}{3}=\dfrac{-d^3+4d}{3}.$$	
	Notice that the lower bound is nothing but the claimed signature $\sigma(X)$ and together with the transversality condition, we may say that $\chi(C_l\cap C_{l+1})$ is always even. Therefore, we have the necessary count. 
	
	In the case  $d=2k+1$, a similar computation applies. The only difference is that one must additionally include the contributions coming from the coamoeba, namely one circle and one isolated point:
	$$\chi(X)=\chi(S^1)+\chi({\text\{point\}})+\sum_{l=0}^{k-1}\big(\chi(\mathbb{RP}^1\times\mathbb{RP}^1)+\chi(C_l)+\chi(C_{l+1})-\chi(C_l\cap C_{l+1})\big),$$
	where $C_l$ is a bidegree $(2l+1,2l+1)$ smooth curve in $\mathbb{RP}^1\times\mathbb{RP}^1$ for each $l$. We may write $0\leq\chi(C_l\cap C_{l+1})\leq 2(2l+1)(2l+3)$. This leaves 
	\begin{align*}
		1 \;\geq\; \chi(X)\geq\;
		1-\sum_{l=0}^{k-1} 2(2l+1)(2l+3)=\frac{-8k^3-12k^2+2k+3}{3}&=-\frac{(2k-1)(2k+1)(2k+3)}{3} \\
		&=-\frac{(d-2)d(d+2)}{3} \\
		&=-\frac{d(d^2-4)}{3}.
	\end{align*}	A similar argument as in the first part shows that we have the required formula again.
	
	The computation of $\chi(\overline{D^{I_{\emptyset}}})$ follows from the fact that $\overline{D^{I_{\emptyset}}}=\emptyset$ in the even degree, while coamoeba effect contributes only a single real plane in the odd degree case.
	
	For the last inequality $d\geq\chi(\overline{D^{I_{S^2}}})\geq d-\sigma(X)$, we may compare it with the previous computations of $\chi(\overline{D^{I_{S^2}}})$.
	
	At each critical level, the contributions of $\chi(\RP^1\times\RP^1)$ are replaced by non-trivial contributions of $\chi(S^2)$. Since $\chi(S^2)=2,$ this produces an additional  $2k$ contribution to $\chi(\overline{D^{I_{S^2}}})$. Moreover, in the odd case, a circle is replaced by a copy of $\RP^2$, contributing an additional  $\chi(\RP^2)=1$. Therefore, in both cases  $d=2k$ or $d=2k+1$, the total additional contribution equals $d$. This establishes the claimed inequalities.
	
\end{proof}

\end{document}
